\sect{कार्त्तिक-कृष्ण-रमा-एकादशी-माहात्म्यम्}
\label{sec:padma-karttika-krishna-rama}

\ganapatyadiDhyanam
\padmaPuranam

\dnsub{अथ द्विषष्टिमोऽध्यायः॥६२}

\uvacha{युधिष्ठिर उवाच}

\twolineshloka
{कथयस्व प्रसादेन मयि स्नेहाज्जनार्दन}
{कार्त्तिकस्यासिते पक्षे किं नामैकादशी भवेत्}% ॥१॥

\uvacha{श्रीकृष्ण उवाच}

\twolineshloka
{श्रूयतां राजशार्दूल कथयामि तवाग्रतः}
{कार्त्तिके कृष्णपक्षे तु रमा नाम सुशोभना}% ॥२॥

\twolineshloka
{एकादशी समाख्याता महापापहरा परा}
{अस्याः प्रसङ्गतो राजन् माहात्म्यं प्रवदामि ते}% ॥३॥

\twolineshloka
{मुचुकुन्द इति ख्यातो बभूव नृपतिः पुरा}
{देवेन्द्रेण समं तस्य मित्रत्वमभवन्नृप}% ॥४॥

\twolineshloka
{यमेन वरुणेनैव कुबेरेणापि सर्वथा}
{विभीषणेन यस्यैव सखित्वमभवन्नृप}% ॥५॥

\twolineshloka
{विष्णुभक्तः सत्यसन्धो बभूव नृपतिः परः}
{तस्यैवं शासतो राजन् राज्यं निहतकण्टकम्}% ॥६॥

\twolineshloka
{बभूव दुहिता गेहे चन्द्रभागा सरिद् वरा}
{शोभनाय च सा दत्ता चन्द्रसेनसुताय वै}% ॥७॥

\twolineshloka
{स कदाचित्समायातः श्वशुरस्य गृहे नृप}
{एकादशीव्रतदिनं समायातं सुपुण्यदम् }% ॥८॥

\twolineshloka
{समागते व्रतदिने चन्द्रभागा त्वचिन्तयत्}
{किं भविष्यति देवेश मम भर्तातिदुर्बलः}% ॥९॥

\twolineshloka
{क्षुधां न क्षमते सोढुं पिता चैवोग्रशासनः}
{पटहस्ताड्यते यस्य सम्प्राप्ते दशमीदिने}% ॥१०॥

\twolineshloka
{न भोक्तव्यं न भोक्तव्यं न भोक्तव्यं हरेर्दिने}
{श्रुत्वा पटहनिर्घोषं शोभनस्त्वब्रवीत् प्रियाम्}% ॥११॥
\onelineshloka*
{किं कर्तव्यं मया कान्ते देहि शिक्षां वरानने}

\uvacha{चन्द्रभागोवाच}
\onelineshloka
{मत्पितुर्वेश्मनि प्रभो भोक्तव्यं नाद्य केनचित्}% ॥१२॥

\twolineshloka
{गजैरश्वैश्च कलभैरन्यैः पशुभिरेव  च}
{तृणमन्नं तथा वारि न भोक्तव्यं हरेर्दिने}% ॥१३॥

\twolineshloka
{मानवैश्च कुतः कान्त भुज्यते हरिवासरे}
{यदि त्वं भोक्ष्यसे कान्त ततो गर्हां प्रयास्यसि}% ॥१४॥
\onelineshloka*
{एवं विचार्य मनसा सुदृढं मानसं कुरु}

\uvacha{शोभन उवाच}
\onelineshloka
{सत्यमेतत् प्रिये वाक्यं करिष्येऽहमुपोषणम्}% ॥१५॥

\twolineshloka
{दैवेन विहितं यद्धि तत्तथैव भविष्यति}
{इति दृढां मतिं कृत्वा चकार नियमं व्रते}% ॥१६॥

\twolineshloka
{क्षुधया पीडिततनुः स बभूवातिदुःखितः}
{इति चिन्तयतस्तस्य आदित्योऽस्तमगाद्गिरिम्}% ॥१७॥

\twolineshloka
{वैष्णवानां नराणां सा निशा हर्षविवर्धनी}
{हरिपूजारतानां च जागरासक्तचेतसाम्}% ॥१८॥

\twolineshloka
{बभूव नृपशार्दूल शोभनस्यातिदुःखदा}
{रवेरुदयवेलायां शोभनः पञ्चतां गतः}% ॥१९॥

\twolineshloka
{दाहयामास राजा तं राजयोग्यैश्च दारुभिः}
{चन्द्रभागा नात्मदेहं ददाह पतिना सह}% ॥२०॥

\twolineshloka
{कृत्वौर्ध्वदैहिकं तस्य तस्थौ जनकवेश्मनि}
{शोभनश्च नृपश्रेष्ठ रमाव्रतप्रभावतः}% ॥२१॥

\twolineshloka
{प्राप्तो देवपुरं रम्यं मन्दराचलसानुनि}
{अनुत्तममनाधृष्यमसङ्ख्येयगुणान्वितम्}% ॥२२॥

\twolineshloka
{हेमस्तम्भमयैसौधै रत्नवैडूर्यमण्डितैः}
{स्फाटिकैर्विविधाकारैर्विचित्रैरुपशोभितैः}% ॥२३॥

\threelineshloka
{सिंहासनसमारूढः सुश्वेतच्छत्रचामरः}
{किरीटकुण्डलयुतो हारकेयूरभूषितः}
{स्तूयमानश्च गन्धर्वैरप्सरोगणसेवितः }% ॥२४॥

\twolineshloka
{शोभनः शोभते यत्र राजराजोपरो यथा}
{सोमशर्मेति विख्यातो मुचुकुन्दपुरेऽभवत्}% ॥२५॥

\twolineshloka
{तीर्थयात्राप्रसङ्गेन भ्रमन्विप्रो ददर्श तम्}
{नृपजामातरं ज्ञात्वा तत् समीपं जगाम सः}% ॥२६॥

\twolineshloka
{शोभनोऽपि तदा ज्ञात्वा सोमशर्माणमागतम्}
{आसनादुत्थितः शीघ्रं नमश्चक्रे द्विजोत्तमम्}% ॥२७॥

\twolineshloka
{चकार कुशलप्रश्नं श्वशुरस्य नृपस्य च}
{कान्तायाश्चन्द्रभागायास्तथैव नगरस्य च}% ॥२८॥

\uvacha{सोमशर्मोवाच}

\twolineshloka
{कुशलं वर्तते राजन् श्वशुरस्य गृहे तव}
{चन्द्रभागा कुशलिनी सर्वतः कुशलं पुरे}% ॥२९॥

\twolineshloka
{स्ववृत्तं कथ्यतां राजन् नाऽऽश्चर्यं विद्यतेऽद्भुतम्}
{पुरं विचित्रं रुचिरं न दृष्टं केनचित्क्वचित्}% ॥३०॥
\onelineshloka*
{एतदाचक्ष्व नृपते कुतः प्राप्तमिदं त्वया}

\uvacha{शोभन उवाच}
\onelineshloka
{कार्तिकस्यासिते पक्षे या नामैकादशी रमा}% ॥३१॥

\twolineshloka
{तामुपोष्य मया प्राप्तं द्विजेन्द्र पुरमध्रुवम्}
{ध्रुवं भवति येनैव तत्कुरुष्व द्विजोत्तमः}% ॥३२॥

\uvacha{द्विज उवाच}

\twolineshloka
{कथमध्रुवमेतद्धि कथं हि भवति ध्रुवम्}
{तत्त्वं कथय राजेन्द्र तत्करिष्यामि नान्यथा}% ॥३३॥

\uvacha{शोभन उवाच}

\twolineshloka
{मयैतद् विहितं विप्र श्रद्धाहीनं व्रतोत्तमम्}
{तेनाहमध्रुवं मन्ये ध्रुवं भवति तच्छृणु}% ॥३४॥

\twolineshloka
{मुचुकुन्दस्य दुहिता चन्द्रभागा सुशोभना}
{तस्यै कथय वृत्तान्तं ध्रुवमेतद् भविष्यति}% ॥३५॥

\uvacha{कृष्ण उवाच}

\twolineshloka
{स श्रुत्वा वचनं तस्य मुचुकुन्दपुरं गतः}
{उवाच सर्वं वृत्तान्त्तं चन्द्रभागाग्रतो द्विजः}% ॥३६॥

\uvacha{सोमशर्मोवाच}

\threelineshloka
{प्रत्यक्षं दयितः कान्तस्तव दृष्टो मया शुभे}
{इन्द्रतुल्यमनाधृष्यं दृष्टं तस्य पुरं मया}
{अध्रुवं तेन तत् प्रोक्तं ध्रुवं भवति तत्कुरु}% ॥३७॥

\uvacha{चन्द्रभागोवाच}

\twolineshloka
{तत्र मां नय विप्रर्षे पतिदर्शनलालसाम्}
{आत्मनो व्रतपुण्येन करिष्यामि पुरं ध्रुवम्}% ॥३८॥

\twolineshloka
{आवयोर्द्विजसंयोगो यथा भवति तत्कुरु}
{प्राप्यते हि महत् पुण्यं कृत्वा योगं वियुक्तयोः}% ॥३९॥

\twolineshloka
{इति श्रुत्वा सह तया सोमशर्मा जगाम ह}
{आश्रमं वामदेवस्य मन्दराचलसन्निधौ}% ॥४०॥

\twolineshloka
{वामदेवोऽशृणोत् सर्वं वृत्तान्तं कथितं तयोः}
{अभ्यषिञ्चच्चन्द्रभागां वेदमन्त्रैरथोज्ज्वलाम्}% ॥४१॥

\twolineshloka
{ऋषिमन्त्रप्रभावेन विष्णुवासरसेवनात्}
{दिव्यदेहा बभूवासौ दिव्यां गतिमवाप ह}% ॥४२॥

\twolineshloka
{पत्युः समीपमगमत् प्रहर्षोत्फुल्ललोचना}
{जहर्ष शोभनोऽतीव दृष्ट्वा कान्तां समागताम्}% ॥४३॥

\twolineshloka
{समाहूय स्वके वामे पार्श्वेतां सन्न्यवेशयत्}
{अथोवाच प्रियं हर्षाच्चन्द्रभागा प्रियं वचः}% ॥४४॥

\twolineshloka
{शृणु कान्त हितं वाक्यं यत् पुण्यं विद्यते मयि}
{अष्टवर्षाधिका जाता यदाहं पितृवेश्मनि}% ॥४५॥

\threelineshloka
{ततः प्रभृति यच्चीर्णं मया चैकादशीव्रतम्}
{यथोक्तविधिसंयुक्तं श्रद्धायुक्तेन चेतसा}
{तेन पुण्यप्रभावेन भविष्यति पुरं ध्रुवम्}% ॥४६॥

\twolineshloka
{सर्वकामसमृद्धं च यावदाभूतसम्प्लवम्}
{एवं सा नृपशार्दूल रमते पतिना सह}% ॥४७॥

\twolineshloka
{दिव्यभोगा दिव्यरूपा दिव्याभरणभूषिता}
{शोभनोऽपि तया सार्धं रमते दिव्यविग्रहः}% ॥४८॥

\twolineshloka
{रमाव्रतप्रभावेन मन्दराचलसानुनि}
{चिन्तामणिसमा ह्येषा कामधेनुसमाथ वा}% ॥४९॥

\twolineshloka
{रमाभिधाना नृपते तवाग्रे कथिता मया}
{तस्या माहात्म्यमनघ श्रुतं सर्वं त्वया नृप}% ॥५०॥

\twolineshloka
{मया तवाग्रे कथितं माहात्म्यं पापनाशनम्}
{एकादशीव्रतानां च पक्षयोरुभयोरपि}% ॥५१॥

\twolineshloka
{यथा कृष्णा तथा शुक्ला विभेदं नैव कारयेत्}
{सेवितैकादशी नॄणां भुक्तिमुक्तिप्रदायिनी}% ॥५२॥

\twolineshloka
{धेनुः श्वेता यथा कृष्णा उभयोः सदृशं पयः}
{तथैव तुल्यफलदं स्मृतमेकादशीद्वयम्}% ॥५३॥

\twolineshloka
{एकादशीव्रतानां यो माहात्म्यं शृणुते नरः}
{सर्वपापविनिर्मुक्तो विष्णुलोके महीयते}% ॥५४॥

॥इति श्रीपाद्मे महापुराणे पञ्चपञ्चाशत्साहस्र्यां संहितायामुत्तरखण्डे उमापतिनारदसंवादान्तर्गतकृष्णयुधिष्ठिरसंवादे कार्त्तिक-कृष्ण-रमा-एकादशी-माहात्म्यं नाम द्विषष्टिमोऽध्यायः॥६२॥

\genMangalaShloka

\hyperref[sec:ekadashi_mahatmyam_padma_puranam]{\closesub}
\clearpage

